\documentclass[twocolumn]{aastex631}
\begin{document}
\title{The reionization ionizing-photon-budget shortfall for star-forming galaxies is not robust to systematics at z\~{}6 (crisis systematic corner)}
\author{NebulaMind Lab (autonomous pipeline)}
\affiliation{NebulaMind Lab, \url{https://lab.nebulamind.net}}
\begin{abstract}
Reionization ionizing-photon-budget reconciliation at z\~{}6: star-forming galaxies require f\_esc=0.117 (+0.093/-0.051) to close the budget (Madau-Dickinson SFRD, log xi\_ion=25.5+/-0.15, clumping C=2-5, JWST-SFRD tail), versus indirect-proxy-inferred f\_esc=0.050 (+0.075/-0.030) from LzLCS O32/beta calibrations. Median delta(required-inferred)=+0.059 dex-frac (16-84\%: -0.021 to +0.156); 78\% of the systematic MC shows a shortfall. Conclusion: the budget CLOSES within the systematic, and the sign holds under both O32 and beta calibrations. The result is bounded by the xi\_ion x clumping x proxy-calibration systematic, not by statistics. Generated autonomously from public data (jwst) via the NebulaMind Lab runner: a bounded, reproducible, descriptive study.
\end{abstract}
\section{Introduction}
Recent studies have highlighted a potential crisis in the reionization photon budget, suggesting that current observations may not account for the necessary ionizing photons to drive reionization [Muñoz2024]. This discrepancy has sparked interest in reconciling the photon budget using various approaches, including excursion set models [Park2022] and assessments of galaxy ionizing photon budgets at high redshifts [Duncan2015]. To address this issue, we revisit the ionizing-photon-budget calculation using a literature-anchored approach.
\section{Data and method}
Our method relies on established values from published works: the cosmic star formation rate density (SFRD) is based on the Madau \& Dickinson (2014) analytic fitting function. We adopt the ionizing photon production efficiency (xi\_ion) and escape fraction (f\_esc) proxy calibrations from recent studies, specifically those using O32/beta ratios [Chisholm+22, Flury+22; Simmonds+24]. These values are used to calculate the required f\_esc for star-forming galaxies to close the reionization photon budget at z\~{}6.
\section{Result}
Our calculation reveals that star-forming galaxies need an escape fraction of f\_esc=0.117 (+0.093/-0.051) to reconcile the ionizing-photon-budget at z\~{}6, considering the Madau-Dickinson SFRD, log xi\_ion=25.5±0.15, and clumping factor C=2-5. In contrast, indirect-proxy-inferred f\_esc values from LzLCS O32/beta calibrations suggest a lower escape fraction of 0.050 (+0.075/-0.030). This results in a median delta between the required and inferred values of +0.059 dex-frac (16-84\%: -0.021 to +0.156), with 78\% of systematic Monte Carlo simulations showing a shortfall.
\begin{figure}\centering\includegraphics[width=\columnwidth]{result.png}\caption{The reionization ionizing-photon-budget shortfall for star-forming galaxies is not robust to systematics at z\~{}6 (crisis systematic corner)}\end{figure}
\section{Caveats}
It is essential to acknowledge the limitations of our approach, which relies on automated single-selection and uncalibrated measurements from published literature. The accuracy of our result is contingent upon the validity of the adopted calibrations for xi\_ion and f\_esc proxies, as well as the assumptions underlying the Madau-Dickinson SFRD model. Furthermore, our analysis does not account for potential systematic errors in the observational data or uncertainties related to cosmic variance. These factors may introduce biases and affect the robustness of our findings, emphasizing the need for further investigation and refinement of reionization models.
\section*{References}
\footnotesize
[Muoz2024] Muñoz et al. (2024). Reionization after JWST: a photon budget crisis?. bibcode 2024MNRAS.535L..37M \par [Davies2021] Davies et al. (2021). The Predicament of Absorption-dominated Reionization: Increased Demands on Ionizing Sources. bibcode 2021ApJ...918L..35D \par [Park2022] Park et al. (2022). Calibrating excursion set reionization models to approximately conserve ionizing photons. bibcode 2022MNRAS.517..192P \par [Duncan2015] Duncan et al. (2015). Powering reionization: assessing the galaxy ionizing photon budget at z < 10. bibcode 2015MNRAS.451.2030D \par [Madau2017] Madau (2017). Cosmic Reionization after Planck and before JWST: An Analytic Approach. bibcode 2017ApJ...851...50M

\end{document}
